\chapter{Elektrische schakelingen en vermogen}\label{ch:g11:circuits-and-power}

Zet een schakelaar om: een lamp gaat branden, doordat een generator vele
kilometers verderop lading door een koperen lus in je muur rondduwt.
Eerdere jaren brachten de schakeling in kaart --- stroomsterkte, \omterm{def:g11:circuits-and-power:voltage}{spanning},
de wet van Ohm. Dit hoofdstuk opent haar opnieuw, met het grootboek van
\cref{ch:g10:energy-conservation}: wat de \omterm{def:g10:orders-of-magnitude:unit}{meter} in rekening brengt,
waarom \omterm{def:g11:lenses-and-eye:lens}{dunne} draden gloeien, waarom het hoogspanningsnet op \qty{400}{kV}
draait.

\section{Lading in beweging}

\begin{definition}[Stroomsterkte en de coulomb]\label{def:g11:circuits-and-power:current}
Een \emph{elektrische stroom}\index{elektrische stroom} is een geordende
stroom van lading (in een metaal: elektronen). Zijn sterkte $I$, in
\emph{ampère}\index{ampère} (\unit{A}), is de lading die per seconde door
een doorsnede van de draad gaat: een constante stroom $I$ die een tijd
$t$ loopt, vervoert de lading
\[ q = I\,t , \]
in \emph{coulomb}\index{coulomb} (\unit{C}), met $\qty{1}{C} =
\qty{1}{A.s}$. De elementaire lading is $e = \qty{1.6e-19}{C}$. De
stroomsterkte wordt gemeten met een ampèremeter, in serie geschakeld.
\end{definition}

\begin{example}[Bliksem tegen telefoon]\label{ex:g11:circuits-and-power:coulomb}
Een blikseminslag voert zo'n $\qty{3.0e4}{A}$ gedurende ongeveer
$\qty{1.0e-4}{s}$: $q = 3.0\times10^{4} \times 10^{-4} = \qty{3.0}{C}$.
Je telefoon, op \qty{2.0}{A}, verplaatst dezelfde lading om de $1.5$
seconde.
\end{example}

\begin{definition}[Spanning en de volt]\label{def:g11:circuits-and-power:voltage}
De \emph{spanning}\index{spanning} $U$ tussen twee punten van een
schakeling, in \emph{volt}\index{volt} (\unit{V}), is de \omterm{def:g10:energy-conservation:energy}{energie} die per
\omterm{def:g10:orders-of-magnitude:unit}{eenheid} lading wordt uitgewisseld bij het afleggen van de weg ertussen:
$\qty{1}{V} = \qty{1}{J/C}$, dus wisselt een lading $q$ die een component
onder spanning $U$ doorkruist er de \omterm{def:g10:energy-conservation:energy}{energie} $E = qU$ mee uit. Een
voltmeter, parallel geschakeld, meet haar.
\end{definition}

\begin{proposition}[Wetten van de schakeling]\label{prop:g11:circuits-and-power:laws}
In een schakeling die in een stationaire toestand werkt, geldt:
\begin{enumerate}
  \item de \emph{knooppuntregel}: de stromen die in een knooppunt
        aankomen tellen op tot de stromen die eruit vertrekken;
  \item de \emph{serieregel}: \omterm{def:g11:circuits-and-power:voltage}{spanningen} langs een weg tellen op:
        $U_{AC} = U_{AB} + U_{BC}$;
  \item de \emph{parallelregel}: twee takken tussen dezelfde twee punten
        staan onder dezelfde \omterm{def:g11:circuits-and-power:voltage}{spanning}.
\end{enumerate}
\end{proposition}

\begin{proof}
In een stationaire toestand hoopt lading zich nergens op, dus stroomt
eruit wat er een knooppunt in stroomt. De \omterm{def:g10:energy-conservation:energy}{energie} per \omterm{def:g11:circuits-and-power:current}{coulomb} telt op
langs een weg en hangt alleen van de eindpunten af --- niet van de tak
die ertussen wordt genomen.
\end{proof}

\section{De wet van Ohm en netwerken van weerstanden}

\begin{definition}[Weerstand]\label{def:g11:circuits-and-power:resistance}
De \emph{weerstand}\index{weerstand} van een geleider is de verhouding
$R = U/I$ van de \omterm{def:g11:circuits-and-power:voltage}{spanning} erover tot de stroomsterkte erdoor, gemeten in
\emph{ohm}\index{ohm} (\unit{\ohm}): $\qty{1}{\ohm} = \qty{1}{V/A}$.
\end{definition}

\begin{proposition}[Wet van Ohm]\label{prop:g11:circuits-and-power:ohm}
\index{wet van Ohm}
Voor een metalen geleider op vaste temperatuur is $R$ constant: \omterm{def:g11:circuits-and-power:voltage}{spanning}
en stroomsterkte zijn evenredig,
\[ U = R\,I . \]
\end{proposition}
\admitted

\begin{remark}\label{rem:g11:circuits-and-power:ohm-limits}
Een experimentele wet, geen universele: een diode, de gloeidraad van een
lamp of je eigen huid gehoorzamen er niet aan. Waarom metalen dat wel
doen, wordt in de universitaire volumes afgeleid.
\end{remark}

\begin{proposition}[Weerstanden in serie en parallel]\label{prop:g11:circuits-and-power:networks}
\index{vervangingsweerstand}
Twee \omterm{def:g11:circuits-and-power:resistance}{weerstanden} in serie zijn gelijkwaardig aan één \omterm{def:g11:circuits-and-power:resistance}{weerstand}
$R_{\text{s}} = R_1 + R_2$; twee parallel aan één \omterm{def:g11:circuits-and-power:resistance}{weerstand}
$R_{\text{p}}$ gegeven door
\[
\frac{1}{R_{\text{p}}} = \frac{1}{R_1} + \frac{1}{R_2} ,
\qquad\text{dus}\quad
R_{\text{p}} = \frac{R_1 R_2}{R_1 + R_2} < \min(R_1, R_2) .
\]
\end{proposition}

\begin{proof}
In serie: dezelfde $I$ doorkruist beide en de \omterm{def:g11:circuits-and-power:voltage}{spanningen} tellen op
(\cref{prop:g11:circuits-and-power:laws}), dus $U = (R_1 + R_2)\,I$.
Parallel: beide staan onder dezelfde $U$ en de stromen tellen op, dus
$I = \frac{U}{R_1} + \frac{U}{R_2}$.
\end{proof}

\begin{omfigure}
\begin{circuitikz}[font=\small]
  \draw (0,0) to[battery1, l=$U$] (0,2.2)
        to[ammeter] (1.9,2.2)
        to[R=$R_1$] (3.7,2.2)
        to[R=$R_2$] (5.5,2.2) -- (5.5,0) -- (0,0);
  \draw (3.7,2.2) -- (3.7,3.3) to[voltmeter] (5.5,3.3) -- (5.5,2.2);
\end{circuitikz}\qquad
\begin{circuitikz}[font=\small]
  \draw (0,0) to[battery1, l=$U$] (0,2.2)
        to[ammeter, l=$I$] (2.2,2.2);
  \draw (2.2,2.2) to[R=$R_1$, i=$I_1$] (2.2,0) -- (0,0);
  \draw (2.2,2.2) -- (4.0,2.2) to[R=$R_2$, i=$I_2$] (4.0,0) -- (2.2,0);
\end{circuitikz}

{\small In serie (links): één stroom, de \omterm{def:g11:circuits-and-power:voltage}{spanningen} tellen op. Parallel
(rechts): één \omterm{def:g11:circuits-and-power:voltage}{spanning}, de stromen tellen op.}
\end{omfigure}

\begin{example}[Vervangingsweerstand]\label{ex:g11:circuits-and-power:networks}
$R_1 = \qty{100}{\ohm}$ en $R_2 = \qty{150}{\ohm}$: in serie
\qty{250}{\ohm}; parallel $\frac{100 \times 150}{250} = \qty{60}{\ohm}$
--- kleiner dan elk van beide: de stroom krijgt een tweede weg.
\end{example}

\section{Energie en vermogen in een schakeling}

\begin{proposition}[Elektrisch vermogen]\label{prop:g11:circuits-and-power:power}
\index{elektrisch vermogen}
Een component onder \omterm{def:g11:circuits-and-power:voltage}{spanning} $U$ waar een stroom $I$ doorheen gaat,
ontvangt \omterm{def:g10:energy-conservation:energy}{energie} met het tempo
\[ P = U\,I , \qquad\text{dus}\quad E = U\,I\,t \ \text{in een tijd } t . \]
\end{proposition}

\begin{proof}
In een tijd $t$ doorkruist de lading $q = It$ de component
(\cref{def:g11:circuits-and-power:current}) en draagt elke \omterm{def:g11:circuits-and-power:current}{coulomb} er $U$
\omterm{def:g10:energy-conservation:energy}{joule} aan over (\cref{def:g11:circuits-and-power:voltage}):
$E = qU = UIt$.
\end{proof}

\begin{proposition}[Vermogen in een weerstand]\label{prop:g11:circuits-and-power:ri2}
Een \omterm{def:g11:circuits-and-power:resistance}{weerstand} $R$ waar een stroom $I$ doorheen gaat, zet
\[ P = R\,I^2 = \frac{U^2}{R} \]
om in warmte.
\end{proposition}

\begin{proof}
Vul de wet van Ohm in $P = UI$ in: $P = (RI)I = RI^2 = U(U/R) = U^2/R$.
\end{proof}

\begin{example}[Een typeplaatje lezen]\label{ex:g11:circuits-and-power:nameplate}
Een waterkoker met het opschrift ``\qty{230}{V} -- \qty{2200}{W}'' trekt
$I = P/U = 2200/230 = \qty{9.6}{A}$ en heeft \omterm{def:g11:circuits-and-power:resistance}{weerstand}
$R = U^2/P = 230^2/2200 = \qty{24}{\ohm}$. In \qty{150}{s} verbruikt hij
$E = 2200 \times 150 = \qty{3.3e5}{J} \approx \qty{0.09}{kW.h}$ ---
ongeveer twee cent.
\end{example}

\begin{remark}[Joules en kilowattuur]\label{rem:g11:circuits-and-power:kwh}
De \omterm{def:g10:orders-of-magnitude:unit}{meter} telt \omterm{def:g10:energy-conservation:kwh}{kilowattuur} (\cref{def:g10:energy-conservation:kwh}):
$\qty{1}{kW.h} = \qty{3.6e6}{J}$. Joules zijn de \omterm{def:g10:orders-of-magnitude:unit}{eenheid} van de
natuurkundige; \omterm{def:g10:energy-conservation:kwh}{kilowattuur} bepalen de prijs op de rekening.
\end{remark}

\begin{omfigure}
\begin{tikzpicture}[font=\small, x=1.5cm, y=0.42cm]
  \draw[-Stealth] (0,0) -- (4.6,0) node[right] {$P$};
  \foreach \x/\t in {1/\qty{10}{W}, 2/\qty{100}{W}, 3/\qty{1}{kW}, 4/\qty{10}{kW}}
    {\draw (\x,-0.3) -- (\x,0.3); \node[below, font=\footnotesize] at (\x,-0.35) {\t};}
  \fill[omDef!70] (0,1) rectangle (1.0,1.75);
  \node[right] at (1.05,1.4) {ledlamp (\qty{10}{W})};
  \fill[omDef!70] (0,2.35) rectangle (1.78,3.1);
  \node[right] at (1.83,2.75) {laptop (\qty{60}{W})};
  \fill[omThm!70] (0,3.7) rectangle (2.95,4.45);
  \node[right] at (3.0,4.1) {broodrooster (\qty{900}{W})};
  \fill[omThm!70] (0,5.05) rectangle (3.34,5.8);
  \node[right] at (3.39,5.45) {waterkoker (\qty{2.2}{kW})};
  \fill[omThm!70] (0,6.4) rectangle (3.95,7.15);
  \node[right] at (4.0,6.8) {elektrische douche (\qty{9}{kW})};
\end{tikzpicture}

{\small \omterm{def:g10:energy-conservation:power}{Vermogens} van apparaten op een logaritmische as: alles wat
verwarmt (rood) leeft een factor tien boven alles wat rekent of schijnt
(blauw).}
\end{omfigure}

\section{Het joule-effect: van broodrooster tot hoogspanningsmast}

\begin{definition}[Joule-effect]\label{def:g11:circuits-and-power:joule}
Het omzetten $P = RI^2$ van \omterm{def:g10:energy-conservation:forms}{elektrische energie} in warmte in elke
\omterm{def:g11:circuits-and-power:resistance}{weerstand} heet het \emph{joule-effect}\index{joule-effect}. Kachels,
waterkokers, broodroosters en ouderwetse gloeilampen zijn met opzet
\omterm{def:g11:circuits-and-power:resistance}{weerstanden}; een \emph{zekering}\index{zekering} is een \omterm{def:g11:lenses-and-eye:lens}{dunne} draad die
zo is afgeregeld dat hij smelt --- en de schakeling verbreekt --- zodra
de stroom zijn nominale waarde overschrijdt.
\end{definition}

\begin{example}[Waarom masten hoog staan]\label{ex:g11:circuits-and-power:transmission}
Een dorp heeft $P = \qty{100}{kW}$ nodig door een lijn met \omterm{def:g11:circuits-and-power:resistance}{weerstand}
$R_\ell = \qty{5.0}{\ohm}$. Geleverd op $U = \qty{500}{V}$:
$I = P/U = \qty{200}{A}$ en de lijn verspilt $R_\ell I^2 = \qty{200}{kW}$
--- twee keer de levering: absurd. Op $U = \qty{20}{kV}$:
$I = \qty{5.0}{A}$ en $R_\ell I^2 = \qty{125}{W}$. Veertig keer minder
stroom, $1600$ keer minder verlies.
\end{example}

\begin{omfigure}
\begin{circuitikz}[font=\small]
  \draw (0,0) to[sV, l=centrale] (0,2)
        to[R, l_=$R_\ell$, i=$I$] (3.4,2) -- (4.8,2)
        to[lamp, l=dorp] (4.8,0) -- (0,0);
  \node[font=\footnotesize, omThm] at (1.7,2.75) {verlies $R_\ell I^2$};
\end{circuitikz}

{\small De levering uit \cref{ex:g11:circuits-and-power:transmission}:
$U$ verhogen bij vaste $P = UI$ krimpt $I$, en het verlies daalt met
$I^2$.}
\end{omfigure}

\begin{remark}[De ijzeren regel van het net]\label{rem:g11:circuits-and-power:grid}
Bij vast geleverd \omterm{def:g10:energy-conservation:power}{vermogen} $P = UI$ daalt het lijnverlies
$R_\ell I^2 = R_\ell P^2/U^2$ met het \emph{kwadraat} van de
transportspanning. Vandaar: omhoog transformeren (\qty{400}{kV}) om te
reizen, omlaag (\qty{230}{V}) om mee te leven.
\end{remark}

\begin{remark}[Veiligheid in huis]\label{rem:g11:circuits-and-power:safety}
Een lijn van \qty{230}{V} met een \omterm{def:g11:circuits-and-power:joule}{zekering} van \qty{16}{A} levert
hoogstens $230 \times 16 \approx \qty{3.7}{kW}$: twee waterkokers
tegelijk laten de \omterm{def:g11:circuits-and-power:joule}{zekering} smelten --- met opzet de dunste, zwakste
schakel, veel goedkoper dan brand in de muur.
\end{remark}

\section{Echte batterijen: elektromotorische spanning}

\begin{definition}[Emk en inwendige weerstand]\label{def:g11:circuits-and-power:emf}
Een batterij zet \omterm{def:g10:energy-conservation:forms}{chemische energie} om in elektrische. Haar
\emph{elektromotorische spanning}\index{elektromotorische spanning} (emk)
$\mathcal{E}$, in \omterm{def:g11:circuits-and-power:voltage}{volt}, is de \omterm{def:g10:energy-conservation:energy}{energie} die ze aan elke \omterm{def:g11:circuits-and-power:current}{coulomb} meegeeft
die haar doorkruist; haar \emph{inwendige weerstand}\index{inwendige
weerstand} $r$ verantwoordt het \omterm{def:g11:circuits-and-power:joule}{joule-effect} in haar eigen chemie.
Model: een ideale bron $\mathcal{E}$ in serie met $r$ (met de sierlijke
$\mathcal{E}$: $E$ is de \omterm{def:g10:energy-conservation:energy}{energie}).
\end{definition}

\begin{proposition}[Klemspanning]\label{prop:g11:circuits-and-power:terminal}
Een batterij die een stroom $I$ levert, toont aan haar klemmen
\[ U = \mathcal{E} - r\,I . \]
\end{proposition}

\begin{proof}
De chemie levert $\mathcal{E}I$, de \omterm{def:g11:circuits-and-power:emf}{inwendige weerstand} eet $rI^2$ op,
de schakeling ontvangt $UI$: het \omterm{thm:g10:energy-conservation:conservation}{behoud van energie}
(\cref{ch:g10:energy-conservation}) geeft $UI = \mathcal{E}I - rI^2$.
\end{proof}

\begin{example}[Een batterij onder belasting]\label{ex:g11:circuits-and-power:discharge}
Een batterij ($\mathcal{E} = \qty{9.0}{V}$, $r = \qty{1.0}{\ohm}$) voedt
een lampje met \omterm{def:g11:circuits-and-power:resistance}{weerstand} $R = \qty{8.0}{\ohm}$. De lus voldoet aan
$\mathcal{E} = rI + RI$, dus $I = \mathcal{E}/(R + r) = \qty{1.0}{A}$ en
$U = 9.0 - 1.0 \times 1.0 = \qty{8.0}{V}$: het lampje ontvangt
\qty{8.0}{W} terwijl de batterij \qty{1.0}{W} verspilt aan zichzelf
opwarmen.
\end{example}

\begin{omfigure}
\begin{tikzpicture}
\begin{axis}[omaxis, width=8cm, height=5cm,
    xlabel={$I$ (\unit{A})}, ylabel={$U$ (\unit{V})},
    xmin=0, xmax=4.6, ymin=0, ymax=10.8,
    xtick={0,1,2,3,4}, ytick={0,3,6,9}]
  \addplot[omThm, thick, domain=0:4.3] {9 - x};
  \addplot[omDef, only marks, mark=*, mark size=1.5pt]
    coordinates {(0.5,8.5) (1.0,8.0) (2.0,7.0) (3.0,6.0) (4.0,5.0)};
  \node[font=\small, anchor=west] at (axis cs:0.2,10.1)
    {snijpunt $\mathcal{E} = \qty{9.0}{V}$};
\end{axis}
\end{tikzpicture}

{\small Ontlaadkarakteristiek van de batterij uit
\cref{ex:g11:circuits-and-power:discharge}: de meetpunten liggen op de
rechte $U = \mathcal{E} - rI$, met richtingscoëfficiënt
$-r = \qty{-1.0}{V/A}$.}
\end{omfigure}

\begin{remark}[Verbruikers]\label{rem:g11:circuits-and-power:receivers}
Een motor of een batterij die wordt opgeladen, laat de omzetting
achterstevoren lopen: hij neemt op onder $U = \mathcal{E}' + r'I$, zet
$\mathcal{E}'I$ om in beweging of chemie en verliest $r'I^2$. Bronnen
duwen, verbruikers duwen terug.
\end{remark}

\begin{example}[Rendement van een batterij]\label{ex:g11:circuits-and-power:efficiency}
Het \omterm{def:g10:energy-conservation:efficiency}{rendement} (\cref{def:g10:energy-conservation:efficiency}) van de
ontlading hierboven: nuttig $UI = \qty{8.0}{W}$ van de
$\mathcal{E}I = \qty{9.0}{W}$, dus $\eta = U/\mathcal{E} = 8/9 \approx
0.89$ --- en het zakt naarmate de stroom stijgt: batterijen die hard
werken, worden heet.
\end{example}

\section{Oefeningen}

\begin{exercise}[$\star$]\label{exo:g11:circuits-and-power:1}
Een telefoonaccu laadt \qty{90}{min} lang met een constante
\qty{2.0}{A}. Welke lading gaat er door de kabel, in \omterm{def:g11:circuits-and-power:current}{coulomb}? Hoeveel
elementaire ladingen zijn dat ($e = \qty{1.6e-19}{C}$)?
\end{exercise}

\begin{exercise}[$\star$]\label{exo:g11:circuits-and-power:2}
Een \omterm{def:g11:circuits-and-power:resistance}{weerstand} van \qty{220}{\ohm} staat onder \qty{12}{V}: welke stroom
loopt er? Een verwarmingselement trekt \qty{0.50}{A} bij \qty{230}{V}:
welke \omterm{def:g11:circuits-and-power:resistance}{weerstand}?
\end{exercise}

\begin{exercise}[$\star$]\label{exo:g11:circuits-and-power:3}
Bereken de vervangingsweerstand van $R_1 = \qty{120}{\ohm}$ en
$R_2 = \qty{60}{\ohm}$ in serie en daarna parallel. Welke is kleiner dan
beide afzonderlijk, en waarom?
\end{exercise}

\begin{exercise}[$\star$]\label{exo:g11:circuits-and-power:4}
Een reisstrijkijzer trekt \qty{8.7}{A} bij \qty{230}{V}. Bereken zijn
\omterm{def:g10:energy-conservation:power}{vermogen} en daarna de \omterm{def:g10:energy-conservation:energy}{energie} die het in \qty{4.0}{min} verbruikt, in
\omterm{def:g10:energy-conservation:energy}{joule} en in \omterm{def:g10:energy-conservation:kwh}{kilowattuur}.
\end{exercise}

\begin{exercise}[$\star$]\label{exo:g11:circuits-and-power:5}
Een laptoplader levert \qty{60}{W} gedurende \qty{5.0}{h} per dag: welke
\omterm{def:g10:energy-conservation:energy}{energie} per dag in \omterm{def:g10:energy-conservation:kwh}{kilowattuur} en in \omterm{def:g10:energy-conservation:energy}{joule}, en welke kosten bij $0.25$
euro per \unit{kW.h}?
\end{exercise}

\begin{exercise}[$\star\star$]\label{exo:g11:circuits-and-power:6}
Een elektrische kachel is aangeduid als \qty{1200}{W} bij \qty{230}{V}.
\begin{enumerate}
  \item Bereken zijn \omterm{def:g11:circuits-and-power:resistance}{weerstand} en de stroom die hij trekt.
  \item Aangesloten (met dezelfde \omterm{def:g11:circuits-and-power:resistance}{weerstand}) op een net van \qty{115}{V}:
        welk \omterm{def:g10:energy-conservation:power}{vermogen} levert hij dan? Waarom niet de helft?
\end{enumerate}
\end{exercise}

\begin{exercise}[$\star\star$]\label{exo:g11:circuits-and-power:7}
$R_1 = \qty{40}{\ohm}$ en $R_2 = \qty{20}{\ohm}$ staan in serie onder
\qty{12}{V}. Bereken de stroom, de \omterm{def:g11:circuits-and-power:voltage}{spanning} over elke \omterm{def:g11:circuits-and-power:resistance}{weerstand} en het
\omterm{def:g10:energy-conservation:power}{vermogen} dat elk omzet; ga na dat de \omterm{def:g10:energy-conservation:power}{vermogens} optellen tot $UI$.
\end{exercise}

\begin{exercise}[$\star\star$]\label{exo:g11:circuits-and-power:8}
$R_1 = \qty{60}{\ohm}$ en $R_2 = \qty{30}{\ohm}$ staan parallel onder
\qty{12}{V}. Bereken de stroom in elke tak, de totale stroom en de
vervangingsweerstand op twee manieren.
\end{exercise}

\begin{exercise}[$\star\star$]\label{exo:g11:circuits-and-power:9}
Een batterij levert $U = \qty{5.5}{V}$ bij $I = \qty{0.50}{A}$ en
$U = \qty{4.0}{V}$ bij $I = \qty{2.0}{A}$. Bepaal haar emk
$\mathcal{E}$ en haar \omterm{def:g11:circuits-and-power:emf}{inwendige weerstand} $r$, en daarna de
kortsluitstroom $\mathcal{E}/r$.
\end{exercise}

\begin{exercise}[$\star\star$]\label{exo:g11:circuits-and-power:10}
Een keukengroep van \qty{230}{V} is beveiligd met een \omterm{def:g11:circuits-and-power:joule}{zekering} van
\qty{16}{A}. Een waterkoker van \qty{2.2}{kW} en een broodrooster van
\qty{1.0}{kW} draaien samen: houdt de \omterm{def:g11:circuits-and-power:joule}{zekering} het? Iemand zet er een
kachel van \qty{1.5}{kW} bij: toon aan dat de \omterm{def:g11:circuits-and-power:joule}{zekering} doorslaat.
\end{exercise}

\begin{exercise}[$\star\star$]\label{exo:g11:circuits-and-power:11}
Een werkplaats heeft \qty{10}{kW} nodig door een voedingskabel met
\omterm{def:g11:circuits-and-power:resistance}{weerstand} \qty{2.0}{\ohm}. Bereken de stroom en het jouleverlies als het
\omterm{def:g10:energy-conservation:power}{vermogen} wordt geleverd op \qty{200}{V}, en daarna op \qty{4.0}{kV}.
Vergelijk de twee verliezen en verklaar de verhouding.
\end{exercise}

\begin{exercise}[$\star\star\star$]\label{exo:g11:circuits-and-power:12}
Noem, met telkens alle drie de gelijke \omterm{def:g11:circuits-and-power:resistance}{weerstanden} van \qty{60}{\ohm},
elke verschillende vervangingsweerstand (serie, parallel, gemengd);
bereken ze alle.
\end{exercise}

\begin{exercise}[$\star\star\star$]\label{exo:g11:circuits-and-power:13}
Een batterij ($\mathcal{E} = \qty{4.5}{V}$, $r = \qty{1.5}{\ohm}$) voedt
een \omterm{def:g11:circuits-and-power:resistance}{weerstand} $R = \qty{6.0}{\ohm}$.
\begin{enumerate}
  \item Bereken $I$, $U$, het nuttige \omterm{def:g10:energy-conservation:power}{vermogen} in $R$, het \omterm{def:g10:energy-conservation:power}{vermogen} dat
        in $r$ verloren gaat, en het \omterm{def:g10:energy-conservation:efficiency}{rendement}.
  \item Toon aan dat het \omterm{def:g10:energy-conservation:power}{vermogen} dat een instelbare $R$ ontvangt,
        $P = \mathcal{E}^2 R/(R+r)^2$, het grootst is wanneer $R = r$
        (aanwijzing: $(R+r)^2/R = (R-r)^2/R + 4r$), en bereken dat
        maximum.
\end{enumerate}
\end{exercise}

\begin{exercise}[$\star\star\star$]\label{exo:g11:circuits-and-power:14}
Een waterkoker van \qty{2.2}{kW} met \omterm{def:g10:energy-conservation:efficiency}{rendement} $0.90$ brengt
\qty{1.0}{kg} water van \qty{20}{\degreeCelsius} aan de kook (water
verwarmen kost \qty{4180}{J} per \omterm{def:g10:orders-of-magnitude:unit}{kilogram} per graad). Bereken de
benodigde warmte, de opgenomen \omterm{def:g10:energy-conservation:forms}{elektrische energie}, de verwarmingstijd en
de kosten bij $0.25$ euro per \unit{kW.h}.
\end{exercise}

\begin{exercise}[$\star\star\star$]\label{exo:g11:circuits-and-power:15}
Een auto-accu: $\mathcal{E} = \qty{12.7}{V}$, $r = \qty{0.020}{\ohm}$,
capaciteit \qty{60}{A.h}.
\begin{enumerate}
  \item Zet de capaciteit om naar \omterm{def:g11:circuits-and-power:current}{coulomb} en schat de opgeslagen \omterm{def:g10:energy-conservation:energy}{energie}
        ($E \approx \mathcal{E} q$) in \omterm{def:g10:energy-conservation:energy}{joule} en in \omterm{def:g10:energy-conservation:kwh}{kilowattuur}.
  \item De startmotor trekt \qty{150}{A}: bereken de klemspanning, het
        geleverde \omterm{def:g10:energy-conservation:power}{vermogen} en het \omterm{def:g10:energy-conservation:power}{vermogen} dat inwendig verloren gaat.
  \item Leg uit waarom de koplampen dimmen terwijl de motor wordt
        aangeslingerd.
\end{enumerate}
\end{exercise}

\section{Probleem: Van de stuwdam tot het broodrooster}

\begin{problem}[{Weekendopgave --- de reis van duizend kilometer van één
kilowattuur, van vallend water tot bruinend brood, met de volledige
rekening van de tocht}]
\label{pb:g11:circuits-and-power:1}
Een stuwdam in de bergen houdt water \qty{200}{m} boven zijn turbines,
die \qty{50}{m^3} --- \qty{5.0e4}{kg} --- per seconde verzwelgen; turbine
en generator zetten om met \omterm{def:g10:energy-conservation:efficiency}{rendement} $0.90$. Door een lijn met totale
\omterm{def:g11:circuits-and-power:resistance}{weerstand} $R_\ell = \qty{32}{\ohm}$ (\qty{1000}{km} heen en terug) voedt
de centrale een stad, waar je broodrooster van \qty{900}{W} op
\qty{230}{V} staat te wachten. De generator levert \qty{20}{kV};
transformatoren (elk met \omterm{def:g10:energy-conservation:efficiency}{rendement} $0.99$) zetten de \omterm{def:g11:circuits-and-power:voltage}{spanning} naar
believen om. Elektriciteit kost $0.25$ euro per \unit{kW.h};
$g = \qty{9.81}{N/kg}$.

\medskip
\noindent\textbf{Deel I --- Bij de stuwdam.}
\begin{enumerate}
  \item Hoeveel \omterm{def:g10:energy-conservation:potential}{potentiële energie} maakt het vallende water elke seconde
        vrij? Leid het \omterm{def:g10:energy-conservation:power}{vermogen} af dat voor de turbines beschikbaar is.
  \item Welk \omterm{prop:g11:circuits-and-power:power}{elektrisch vermogen} verlaat de generator?
  \item Bij \qty{20}{kV} rechtstreeks van de generator: welke stroom
        loopt er?
  \item Bereken het jouleverlies $R_\ell I^2$ bij die stroom en vergelijk
        het met wat de centrale levert. Trek je conclusie.
  \item Het \omterm{def:g10:energy-conservation:power}{vermogen} van de stad $P = UI$ ligt vast door de vraag. Wat is
        de enige manier om $I$ te verkleinen, en welk apparaat doet dat?
\end{enumerate}

\medskip
\noindent\textbf{Deel II --- Op de lijn.}
\begin{enumerate}[resume]
  \item Een transformator, voor deze vraag verliesvrij, brengt de
        \qty{88}{MW} omhoog naar \qty{400}{kV}: welke stroom loopt er nu?
  \item Bereken het nieuwe lijnverlies, in megawatt.
  \item Welk deel van wat de centrale levert, gaat in de lijn verloren?
  \item Toon aan dat het verloren deel gelijk is aan $R_\ell P/U^2$, en
        ga na dat dit de factor tussen vraag 4 en vraag 7 verklaart.
  \item Bereken het spanningsverlies $R_\ell I$ over de lijn bij
        \qty{400}{kV}. Klopt dat met vraag 8?
\end{enumerate}

\medskip
\noindent\textbf{Deel III --- In de keuken.}
\begin{enumerate}[resume]
  \item Zet \qty{1}{kW.h} om naar \omterm{def:g10:energy-conservation:energy}{joule}; wat vermenigvuldigt de \omterm{def:g10:orders-of-magnitude:unit}{meter}?
  \item Bereken de stroom en de \omterm{def:g11:circuits-and-power:resistance}{weerstand} van het broodrooster bij
        \qty{230}{V}.
  \item Drie minuten roosteren: de \omterm{def:g10:energy-conservation:energy}{energie} in \omterm{def:g10:energy-conservation:energy}{joule} en in \omterm{def:g10:energy-conservation:kwh}{kilowattuur}, en
        de prijs ervan.
  \item Het broodrooster zet $RI^2$ om: leg in één zin uit waarom
        dezelfde formule op de lijn het verlies is en in de keuken het
        product.
  \item Het broodrooster deelt een \omterm{def:g11:circuits-and-power:joule}{zekering} van \qty{16}{A} met een
        waterkoker van \qty{2.2}{kW}: houdt de \omterm{def:g11:circuits-and-power:joule}{zekering} het? Er komt een
        kachel van \qty{2.0}{kW} bij: en nu?
\end{enumerate}

\medskip
\noindent\textbf{Deel IV --- De doorlichting.}
\begin{enumerate}[resume]
  \item De huisbedrading zelf, met \omterm{def:g11:circuits-and-power:resistance}{weerstand} \qty{0.05}{\ohm}, voert de
        \qty{13.5}{A} uit vraag 15: bereken het verlies en welk deel dat
        is van de ruwweg \qty{3.1}{kW} die wordt geleverd.
  \item Rijg de \omterm{def:g10:energy-conservation:efficiency}{rendementen} aaneen (turbine--generator, optransformeren,
        lijn, aftransformeren, huisbedrading) tot één totaalrendement van
        het net.
  \item Hoeveel \omterm{def:g10:orders-of-magnitude:unit}{kilogram} (en hoeveel kubieke \omterm{def:g10:orders-of-magnitude:unit}{meter}) stuwwater moet er
        vallen om \emph{één} \omterm{def:g10:energy-conservation:kwh}{kilowattuur} bij je broodrooster te krijgen?
  \item Bereken het verloren deel $R_\ell P/U^2$ uit vraag 9 bij
        $U = \qty{20}{kV}$. Wat betekent een waarde boven $1$?
  \item De clou --- \emph{de prijs van de afstand}: zeg in twee zinnen
        wat één geroosterd \omterm{def:g10:energy-conservation:kwh}{kilowattuur} de berg kost, welk deel onderweg
        sneuvelt, en welke ene truc de reis mogelijk maakt.
\end{enumerate}
\end{problem}
